\naslov{Pridobivanje cenilk}

\podnaslov{Metoda empirične porazdelitve}

Opazimo vrednosti $X_1, \ldots, X_n$, ki so enako porazdeljene.
Te predstavljajo vzorec preučevane porazdelitve.
Naj se $y$ izraža kot karakteristika te porazdelitve, $y = g(\theta) =
\operatorname{Char}_\theta(X_1)$, kjer je $\operatorname{Char}_\theta$ odvisna
le od porazdelitve $X_1$ pri verjetnostni meri $P_\theta$.
Cenilka za $y$ je potem $\hat{y} = \operatorname{Char}(\operatorname{Emp}(X_1,
\ldots, X_n))$.

\vprasanje{Opiši metodo empirične porazdelitve. Kakšne so predpostavke?}

Kot primer vzemimo $\mu = E(X_1)$, ki ga ocenimo z $\hat{\mu} =
E(\operatorname{Emp}(X_1, \ldots, X_n)) = \overline{X}_n$.
To je nepristranska cenilka za $\mu$.
Če so $X_1, \ldots, X_n$ nekorelirane, vemo, da je
\[
  \operatorname{SE} = \frac{\sigma}{\sqrt{n}}
\]
kjer je $\sigma^2 = \var(X_1)$.
Drug pogled je enostavno slučajno vzorčenje.
Na populaciji z $N$ enotami naj bo definirana statistična spremenljivka,
tj.~funkcija $\{1, \ldots, N\} \to \R$.
Njene vrednosti na posameznih enotah izrazimo z $x_1, \ldots, x_N$.
Iz populacije vzamemo vzorec, ki ga sestavljajo enote $K_1, \ldots, K_n$.
Enostavno slučajno vzorčenje pomeni, da ima vektor $(K_1, \ldots, K_n)$
vrednosti v množici $\{(k_1, \ldots, k_n) \such k_i \ne k_j\}$.
Označimo $X_i = x_{K_i}$, da velja $X_i \sim \operatorname{Emp}(X_1, \ldots,
X_n)$.
Potem je $\mu = E(X_1) = \mu$ in
\[
  \var(X_1) = \inv{N} \sum_{k=1}^N (x_k - \mu)^2.
\]
Standardna napaka cenilke $\overline{X}$ za $\overline{x}$ je
\[
  \operatorname{SE}^2
  = \var(\overline{X})
  = \inv{n^2} \var\left(\sum_i X_i\right)
  = \inv{n^2} \sum_{i,j} \cov(X_i, X_j)
  = \frac{N-n}{N-1} \frac{\sigma^2}{n}.
\]

\vprasanje{Opiši uporabo metode empirične porazdelitve za pričakovano vrednost.}

Če ocenjujemo $\sigma^2 = \var(X_1)$, je cenilka
\[
  \hat{\sigma}^2 = \inv{n} \sum_k (X_k - \overline{X})^2
\]
Za nepristranskost izračunajmo $E(\hat{\sigma}^2)$.
Velja
\[
  \hat{\sigma}^2 = \inv{n} \norm{(X_1 - \overline{X}, \ldots, X_n -
	\overline{X})}^2
  = \inv{n} \norm{\sv{X} - \overline{X} \cdot \sv{1}}^2
  = \inv{n} \norm{(I - \inv{n} \sv{1} \sv{1}^T) \sv{X}}^2.
\]
Označimo $H = I - \inv{n} \sv{1} \sv{1}^T$.
To je ortogonalni projektor, velja
\[
  \hat{\sigma}^2 = \inv{n} (H \sv{X})^T (H \sv{X})
  = \inv{n} \sv{X}^T H \sv{X}.
\]
Sled skalarja je enaka skalarju, velja
\[
  E(\hat{\sigma}^2)
  = \inv{n} \sled (H E(\sv{X} \sv{X}^T)).
\]
Če zamenjamo spremenljivke $\sv{X}' = \sv{X} - \sv{\mu}$, lahko izrazimo
\[
  E(\hat{\sigma}^2) = \inv{n} \sled(H E(\sv{X}' \sv{X}'^T))
  = \inv{n} \sled (H \var(\sv{X})),
\]
ker je $E(\sv{X}') = 0$.

Sedaj ločimo dva primera.
Če so $X_1, \ldots, X_n$ nekorelirane, je $\var(\sv{X}) = \sigma^2 I$ in
posledično
\[
  E(\hat{\sigma}^2) = \frac{n-1}{n} \sigma^2,
\]
torej je $\hat{\sigma}^2$ pristranska cenilka za $\sigma^2$, cenilka
$\hat{\sigma}^2_+ = \frac{n}{n-1} \hat{\sigma}^2$ pa je nepristranska.

\vprasanje{Pokaži, da je empirična varianca pristranska.}